Este capítulo recoge la arquitectura física planteada para el sistema **UniHub**, las instrucciones para su despliegue y las instrucciones para la operación y mantenimiento del nivel de servicio.

\section{Arquitectura Física}
La arquitectura física del sistema se basa en la virtualización liviana mediante contenedores Docker sobre un servidor anfitrión con sistema operativo Linux o Windows con soporte para Docker Engine. La infraestructura se compone de 4 contenedores aislados que se comunican en una red puente virtualizada (\texttt{ruct\_network}):
\begin{enumerate}
    \item \textbf{Servidor de Base de Datos (\texttt{ruct\_db}):} Contenedor PostgreSQL 15 escuchando en el puerto 5432 con volumen nominado \texttt{ruct\_postgres\_data}.
    \item \textbf{Servidor de API REST (\texttt{ruct\_api}):} Contenedor Python 3.12 / FastAPI escuchando en el puerto 8000.
    \item \textbf{Servidor Web Nginx (\texttt{ruct\_www}):} Contenedor Nginx escuchando en los puertos 80 y 5173.
    \item \textbf{Contenedor de Ingesta (\texttt{ruct\_crawler}):} Contenedor Python para ejecuciones del rastreador con montaje directo en la carpeta del host \texttt{Codigo/Crawler/Datos}.
\end{enumerate}

\section{Instrucciones de despliegue}

\subsection{Requisitos previos}
\begin{itemize}
    \item Docker Engine 20.10+ y Docker Compose v2.
    \item Al menos 2 GB de memoria RAM libre y 5 GB de espacio en disco.
\end{itemize}

\subsection{Inventario de componentes}
Los artefactos del despliegue se ubican en \texttt{Codigo/Docker/}:
\begin{itemize}
    \item \texttt{docker-compose.yml}
    \item \texttt{crawler/Dockerfile} y \texttt{crawler/entrypoint.sh}
    \item \texttt{api/Dockerfile} y \texttt{api/entrypoint.sh}
    \item \texttt{www/Dockerfile} y \texttt{www/nginx.conf}
\end{itemize}

\subsection{Procedimientos de instalación}
Para desplegar todo el sistema contenerizado:
\begin{enumerate}
    \item Navegar a la carpeta del entorno Docker:
        \begin{lstlisting}[language=bash]
cd d:\Proyecto\Codigo\Docker
        \end{lstlisting}
    \item Construir y levantar los servicios en segundo plano:
        \begin{lstlisting}[language=bash]
docker compose up --build -d
        \end{lstlisting}
    \item Poblar la base de datos PostgreSQL desde los archivos JSON:
        \begin{lstlisting}[language=bash]
docker compose exec api python database/etl_loader.py
        \end{lstlisting}
\end{enumerate}

\subsection{Pruebas de implantación}
Acceder desde el navegador web a las siguientes direcciones de verificación:
\begin{itemize}
    \item Portal Web Frontend UniHub: \url{http://localhost} o \url{http://localhost:5173}
    \item Documentación API Swagger: \url{http://localhost:8000/docs}
    \item Panel de Administración (Acceso Aislado): \url{http://localhost/admin}
\end{itemize}

\section{Instrucciones para la operación del sistema y mantenimiento del nivel de servicio}
Para garantizar la continuidad de servicio y la persistencia de los datos:
\begin{itemize}
    \item \textbf{Chequeo de Logs:} Monitorear la salud de los servicios con \texttt{docker compose logs -f}.
    \item \textbf{Copia de Seguridad de la BD:} Exportar backup mediante \texttt{docker compose exec db pg\_dump -U postgres ruct\_db > backup.sql}.
    \item \textbf{Copia de Seguridad de Archivos JSON:} La carpeta \texttt{Codigo/Crawler/Datos} en el host mantiene los archivos descargados e insustituibles físicamente en disco.
\end{itemize}